% gpu_reliability_bars.tex — GPU 可靠性对比柱状图
\begin{tikzpicture}
\begin{axis}[
    width=0.95\textwidth,
    height=6.5cm,
    ybar,
    bar width=28pt,
    ylabel={年化失效率 AFR（\%）},
    ymin=0, ymax=38,
    ytick={0,3,5,10,15,30},
    yticklabels={0,3,5,10,15,30},
    symbolic x coords={地面数据中心, 轨道商业级GPU, 轨道抗辐射GPU, 极端太阳事件},
    xtick=data,
    xticklabel style={align=center, font=\footnotesize},
    enlarge x limits=0.25,
    nodes near coords,
    nodes near coords style={font=\footnotesize},
    point meta=explicit symbolic,
    every node near coord/.append style={yshift=6pt},
    legend style={
        at={(0.5,-0.22)},
        anchor=north,
        legend columns=1,
        draw=none,
        font=\footnotesize
    },
    grid=major,
    grid style={draw=gray!20},
    axis lines=left,
    enlarge y limits={upper, value=0.2}
]

% 柱状图数据
\addplot[fill=primary!60, draw=primary]
    coordinates {
        (地面数据中心, 2) [2\%]
        (轨道商业级GPU, 12) [12\%]
        (轨道抗辐射GPU, 5) [5\%]
        (极端太阳事件, 30) [30\%+]
    };

% 标注
\draw[->, >=stealth, highlight, line width=1.0pt, dashed]
    (axis cs:地面数据中心,8) -- (axis cs:轨道商业级GPU,8)
    node[midway, above, font=\footnotesize, text=highlight] {$\times$6};

\draw[->, >=stealth, highlight, line width=1.0pt, dashed]
    (axis cs:轨道商业级GPU,20) -- (axis cs:极端太阳事件,20)
    node[midway, above, font=\footnotesize, text=highlight] {$\times$2.5};

% 不可接受阈值线
\draw[red!50, dashed, line width=0.8pt]
    (axis cs:地面数据中心,15) -- (axis cs:极端太阳事件,15);
\node[red!60, font=\footnotesize] at (axis cs:轨道抗辐射GPU,16.5)
    {不可接受阈值};

% 注解
\node[align=left, font=\footnotesize, text=gray]
    at (axis cs:地面数据中心,36)
    {在轨维修 $\rightarrow$ 不可行\\集群3年可用算力 $\rightarrow$ 20--40\%};

\end{axis}
\end{tikzpicture}
